\subsubsection{Hopcroft-Karp}
\begin{lstlisting}[style=mystyle, language=C++]
// TC: O(sqrt(V) * E) donde V = nodos lado izquierdo, E = aristas
// usa: matching maximo en grafo bipartito; usa BFS + DFS por capas para encontrar varios augmenting paths por fase
// nota: para N, M <= 1e5; para N, M <= 1e3 alcanza Kuhn (O(V*E))
struct HopcroftKarp {
	int n, m; // n->#left_nodes  m->#right_nodes
	vector<vector<int>> adj;
	vector<int> dist, matchU, matchV;

	HopcroftKarp(int n, int m): n(n), m(m) {
		adj.assign(n + 1, {});
		matchU.assign(n + 1, 0);
		matchV.assign(m + 1, 0);
		dist.assign(n + 1, 0);
	}

	void addEdge(int u, int v) {
		adj[u].push_back(v);
	}

	bool bfs() {
		queue<int> q;
		for (int u = 1; u <= n; u++) {
			if (matchU[u] == 0) {
				dist[u] = 0;
				q.push(u);
			} else {
				dist[u] = -1;
			}
		}
		bool found = false;
		while (!q.empty()) {
			int u = q.front(); q.pop();
			for (int v : adj[u]) {
				if (matchV[v] != 0 && dist[matchV[v]] == -1) {
					dist[matchV[v]] = dist[u] + 1;
					q.push(matchV[v]);
				}
				if (matchV[v] == 0) {
					found = true;
				}
			}
		}
		return found;
	}

	bool dfs(int u) {
		for (int v : adj[u]) {
			if (matchV[v] == 0 ||
			(dist[matchV[v]] == dist[u] + 1 && dfs(matchV[v]))) {
				matchU[u] = v;
				matchV[v] = u;
				return true;
			}
		}
		dist[u] = -1;
		return false;
	}

	int maxMatching() {
		int matching = 0;
		while (bfs()) {
			for (int u = 1; u <= n; u++) {
				if (matchU[u] == 0 && dfs(u)) {
					matching++;
				}
			}
		}
		return matching;
	}
};

int main(){
	// entrada: n nodos izquierda, m nodos derecha, k aristas (u, v) con u en [1..n], v en [1..m]
	int n, m, k;
	cin >> n >> m >> k;
	HopcroftKarp hk(n, m);
	for(int i=0;i<k;i++){
		int u, v; cin >> u >> v;
		hk.addEdge(u, v);
	}
	cout << hk.maxMatching() << "\n";
	return 0;
}
\end{lstlisting}
